\documentclass[12pt]{article}

\usepackage{amsmath, amssymb}
\usepackage{geometry}
\usepackage{titlesec}

\geometry{margin=1in}

\title{CSE400 -- Fundamentals of Probability in Computing\\
Lecture 11: Transformation of Random Variables}
\date{}
\begin{document}

\maketitle

\section*{Lecture Outline}

The lecture covers the following topics:

\begin{enumerate}
\item \textbf{Transformation of Random Variables} \\
Learning transformation techniques for random variables.

\item \textbf{Function of Two Random Variables} \\
Joint transformations and derived distributions.

\item \textbf{Illustrative Example} \\
Detailed derivation for the case
\[
Z = X + Y
\]
\end{enumerate}

These topics together explain how the distribution of a \textbf{new random variable} can be derived from existing random variables.

\section{Transformation of Random Variables}

\subsection*{Definition}

A \textbf{transformation of a random variable} occurs when a new random variable is defined as a function of another random variable.

Let $X$ be a random variable and let

\[
Y = g(X)
\]

Then $Y$ is called a \textbf{transformed random variable}.

The main objective is to determine the \textbf{probability distribution of $Y$} given the probability distribution of $X$.

\section{Distribution of a Transformed Variable}

The most common method for finding the distribution of a transformed variable is the \textbf{CDF method}.

\subsection*{Step 1: Start with the CDF Definition}

The cumulative distribution function of $Y$ is

\[
F_Y(y) = P(Y \leq y)
\]

\subsection*{Step 2: Substitute the Transformation}

If

\[
Y = g(X)
\]

then

\[
F_Y(y) = P(g(X) \leq y)
\]

\subsection*{Step 3: Express Probability in Terms of $X$}

The evaluation of this probability depends on whether the function $g(x)$ is:

\begin{itemize}
\item \textbf{Monotonic increasing}, or
\item \textbf{Monotonic decreasing}
\end{itemize}

\section{Case 1: Monotonic Increasing Transformation}

Assume that the function $g(x)$ is \textbf{strictly increasing}.

Then

\[
g(X) \le y
\]

implies

\[
X \le g^{-1}(y)
\]

Thus,

\[
F_Y(y) = P(X \le g^{-1}(y))
\]

\[
F_Y(y) = F_X(g^{-1}(y))
\]

\subsection*{Derivation of the PDF}

If the random variable $X$ has probability density function $f_X(x)$, then the PDF of $Y$ is obtained by differentiation.

\[
f_Y(y) = \frac{d}{dy}F_Y(y)
\]

Substituting the expression of $F_Y(y)$:

\[
f_Y(y) = \frac{d}{dy}F_X(g^{-1}(y))
\]

Using the \textbf{chain rule}

\[
f_Y(y) =
f_X(g^{-1}(y))
\cdot
\frac{d}{dy}g^{-1}(y)
\]

Thus the transformation formula becomes

\[
f_Y(y) =
f_X(g^{-1}(y))
\left|
\frac{d}{dy}g^{-1}(y)
\right|
\]

\section{Case 2: Monotonic Decreasing Transformation}

If $g(x)$ is \textbf{strictly decreasing}, then

\[
g(X) \le y
\]

implies

\[
X \ge g^{-1}(y)
\]

Therefore

\[
F_Y(y) = P(X \ge g^{-1}(y))
\]

\[
F_Y(y) = 1 - F_X(g^{-1}(y))
\]

The probability density function can again be obtained by differentiating this CDF.

\section{Function of Two Random Variables}

In many situations a new random variable is defined as a \textbf{function of two random variables}.

Let

\[
Z = h(X,Y)
\]

To determine the distribution of $Z$, we require the \textbf{joint probability distribution} of $X$ and $Y$.

\subsection*{Joint Probability Density Function}

If $X$ and $Y$ are continuous random variables, they may have a joint PDF

\[
f_{X,Y}(x,y)
\]

This function describes the probability behavior of the pair $(X,Y)$.

\subsection*{Independence Assumption}

If $X$ and $Y$ are \textbf{independent}, then their joint PDF is

\[
f_{X,Y}(x,y) = f_X(x)f_Y(y)
\]

This property simplifies the derivation of distributions for functions such as

\begin{itemize}
\item $Z = X + Y$
\item $Z = X - Y$
\item $Z = XY$
\end{itemize}

\section{Illustrative Example}

\subsection*{Problem Statement}

Find the distribution of

\[
Z = X + Y
\]

where

\begin{itemize}
\item $X$ and $Y$ are \textbf{independent random variables}
\item Both follow \textbf{exponential distributions}
\end{itemize}

\[
X \sim Exp(\lambda)
\]

\[
Y \sim Exp(\lambda)
\]

\subsection*{Step 1: Write the PDFs}

For an exponential random variable

\[
f_X(x) = \lambda e^{-\lambda x}, \quad x \ge 0
\]

\[
f_Y(y) = \lambda e^{-\lambda y}, \quad y \ge 0
\]

\subsection*{Step 2: Apply the Convolution Formula}

For independent variables

\[
Z = X + Y
\]

the probability density function is

\[
f_Z(z) =
\int_{-\infty}^{\infty}
f_X(x)f_Y(z-x)\,dx
\]

\subsection*{Step 3: Determine the Limits of Integration}

Because exponential variables are defined only for

\[
x \ge 0
\]

and

\[
y = z-x \ge 0
\]

we obtain

\[
0 \le x \le z
\]

Therefore

\[
f_Z(z) =
\int_0^z
f_X(x)f_Y(z-x)\,dx
\]

\subsection*{Step 4: Substitute the PDFs}

\[
f_Z(z) =
\int_0^z
(\lambda e^{-\lambda x})
(\lambda e^{-\lambda(z-x)})
dx
\]

\subsection*{Step 5: Simplify the Expression}

\[
f_Z(z) =
\lambda^2 e^{-\lambda z}
\int_0^z dx
\]

\subsection*{Step 6: Evaluate the Integral}

\[
\int_0^z dx = z
\]

Thus

\[
f_Z(z) =
\lambda^2 z e^{-\lambda z}
\]

\subsection*{Final Probability Density Function}

\[
f_Z(z)=
\begin{cases}
\lambda^2 z e^{-\lambda z}, & z \ge 0 \\
0, & \text{otherwise}
\end{cases}
\]

\subsection*{Interpretation}

The sum of two independent exponential random variables follows a \textbf{Gamma (Erlang) distribution} with:

\begin{itemize}
\item shape parameter $k = 2$
\item rate parameter $\lambda$
\end{itemize}

\section*{Key Formulas from Lecture 11}

\subsection*{Transformation Formula}

\[
f_Y(y) =
f_X(g^{-1}(y))
\left|
\frac{d}{dy} g^{-1}(y)
\right|
\]

\subsection*{Convolution Formula}

\[
f_Z(z) =
\int_{-\infty}^{\infty}
f_X(x)f_Y(z-x)dx
\]

For non-negative variables:

\[
f_Z(z) =
\int_0^z
f_X(x)f_Y(z-x)dx
\]

\end{document}